\documentclass[oneside,11pt,dvipsnames]{book}
\usepackage[dvipsnames]{xcolor}
\usepackage{pagecolor}
\usepackage[utf8]{inputenc}
\usepackage[T1]{fontenc}
\usepackage{soul}
\usepackage[margin=1.5in]{geometry}
\usepackage{tcolorbox}
\tcbuselibrary{documentation, skins}
\usepackage{titlesec}


% \setsecnumdepth{subsection}
% \setcounter{tocdepth}{3}
\usepackage{fontawesome5}
\usepackage{censor}
\usepackage{enumitem}
\usepackage{cite}
\usepackage{caption}
\captionsetup{font=small}
\usepackage{graphicx}
\usepackage{pdfpages}
\usepackage{hyperref}
\usepackage{wrapfig}
\setlength\intextsep{0pt} % remove extra space above and below in-line float

\usepackage{booktabs}

\usepackage{tikz}
\usetikzlibrary{calc}

%\usepackage{anyfontsize}
%\usepackage{sectsty}
\usepackage[acronym]{glossaries}
\setacronymstyle{long-short-desc}
\makeglossaries{}
\usepackage[makeroom]{cancel}

% boxes
\newtcolorbox{marker}[1][]{enhanced,
  before skip balanced=2mm,after skip balanced=3mm,
  boxrule=0.4pt,left=5mm,right=2mm,top=1mm,bottom=1mm,
  colback=yellow!10,
  sharp corners,rounded corners=southeast,arc is angular,arc=3mm,
  underlay={%
    \path[fill=tcbcolback!80!black] ([yshift=3mm]interior.south east)--++(-0.4,-0.1)--++(0.1,-0.2);
    \path[draw=tcbcolframe,shorten <=-0.05mm,shorten >=-0.05mm] ([yshift=3mm]interior.south east)--++(-0.4,-0.1)--++(0.1,-0.2);
    \path[fill=tcbcolback!80!black,draw=none] (interior.south west) rectangle node{\large\bfseries {#1}} ([xshift=4mm]interior.north west);
    },
  drop fuzzy shadow}

\newenvironment{commentbox}[1][]{
  \small
  \begin{marker}[\faChessPawn{}]
    {\textbf{#1}}
  }{
  \end{marker}
}

\newenvironment{surprisebox}[1][]{
  \small
  \begin{marker}[\faChessKnight{}]
    {\textbf{#1}}
  }{
  \end{marker}
}

\newenvironment{examplebox}[1][]{
  \small
  \begin{marker}[\faChessBishop{}]
    {\textbf{#1}}
  }{
  \end{marker}
}

\newenvironment{questionbox}[1][]{
  \small
  \begin{marker}[\faChessRook{}]
    {\textbf{#1}}
  }{
  \end{marker}
}

\newenvironment{keybox}[1][]{
\small
  \begin{marker}[\faChessQueen{}]
    {\textbf{#1}}
  }{
  \end{marker}
}

\newenvironment{warningbox}[1][]{
  \small
  \begin{marker}[\faChessKing{}]
    {\textbf{#1}}
  }{
  \end{marker}
}

\renewcommand{\figurename}{Fig.}
\renewcommand{\tablename}{Tab.}
%\def\Section{\S}
\newcommand{\Section}{\textsection}

\renewcommand{\figureautorefname}{Fig.}
\renewcommand{\tableautorefname}{Tab.}
\makeatletter
% \renewcommand{\chapterautorefname}{\S\@gobble}
% \renewcommand{\sectionautorefname}{\S\@gobble}
% \renewcommand{\subsectionautorefname}{\S\@gobble}
% \renewcommand{\appendixautorefname}{\S\@gobble}
\renewcommand{\chapterautorefname}{\textsection\@gobble}
\renewcommand{\sectionautorefname}{\textsection\@gobble}
\renewcommand{\subsectionautorefname}{\textsection\@gobble}
\renewcommand{\appendixautorefname}{\textsection\@gobble}

\newcommand*\l@chapterinfo{\@nodottedtocline{0}{0.0em}{1.5em}}
\newcommand*\l@sectioninfo{\@nodottedtocline{1}{1.5em}{2.3em}}
\newcommand*\l@subsectioninfo{\@nodottedtocline{2}{3.8em}{3.2em}}
%\newcommand*\l@subsubsectioninfo{\@nodottedtocline{3}{7.0em}{4.1em}}
%\newcommand*\l@paragraphinfo{\@nodottedtocline{4}{10em}{5em}}
%\newcommand*\l@subparagraphinfo{\@nodottedtocline{5}{12em}{6em}}line{5}{12em}{6em}}

% Copied from the book class macro `\@dottedtocline`. Removed the dots and page number
\def\@nodottedtocline#1#2#3#4#5{%
  \ifnum #1>\c@tocdepth \else
    \vskip \z@ \@plus.2\p@
    {\leftskip #2\relax \rightskip \@tocrmarg \parfillskip -\rightskip
     \parindent #2\relax\@afterindenttrue
     \interlinepenalty\@M
     \leavevmode
     \@tempdima #3\relax
     \advance\leftskip \@tempdima \null\nobreak\hskip -\leftskip
     {#4}\nobreak
     \leaders\hbox{$\m@th
        \mkern \@dotsep mu\hbox{\,}\mkern \@dotsep
        mu$}\hfill
     \nobreak
     \hb@xt@\@pnumwidth{\hfil\normalfont \normalcolor }%
     \par}%
  \fi}



\makeatother

\newcommand{\chapterinfo}[1]{%
  \phantomsection
  \addcontentsline{toc}{chapterinfo}{\textcolor{black}{\emph{#1}}}%
}
\newcommand{\sectioninfo}[1]{%
  \phantomsection
  \addcontentsline{toc}{sectioninfo}{\textcolor{black}{\emph{#1}}}%
}
\newcommand{\subsectioninfo}[1]{%
  \phantomsection
  \addcontentsline{toc}{subsectioninfo}{\textcolor{black}{\texttt #1}}%
}

\newcommand{\mycomment}[3][\color{blue}]{{#1{{#2}: {#3}}}}
\newcommand{\tvn}[1]{\mycomment{TVN}{#1}}{}
\newcommand{\didi}[1]{\mycomment{Didier}{#1}}{}
\newcommand{\tl}[1]{\mycomment{ThanhLe}{#1}}{}
\newcommand{\xz}[1]{\mycomment{Xiaokuan}{[#1]}}{}
\newcommand{\red}[1]{{\color{red}{#1}}}

\newcommand{\mytitle}{GMU CS Notes}
\newcommand{\mysubtitle}{Stuff useful for me}
\newcommand{\mybookaddress}{\href{https://roars.dev/phd-cs-us/demystify.pdf}{https://roars.dev/phd-cs-us/demystify.pdf}}
\newcommand{\mybookgithub}{https://code.roars.dev/phd-cs-us}
%\newcommand{\me}{ThanhVu Nguyen}
%\newcommand{\institute}{George Mason University}

\usepackage{hyperref}
\hypersetup{
    hidelinks,
    linktoc=all,
  colorlinks,
  citecolor=black,
  filecolor=black,
  linkcolor=blue,
  urlcolor=blue,
  pdfpagemode=UseOutlines,
}




\begin{document}

\frontmatter
%\newpagecolor{SpringGreen}

\begin{titlepage}
    \begin{tikzpicture}[remember picture, overlay]
        \draw[OliveGreen, line width = 3pt] ($(current page.north west) + (1in,-1in)$) rectangle ($(current page.south east) + (-1in,1in)$);
      \end{tikzpicture}

    \centering
    {\Huge\bfseries\baselineskip1.3em\mytitle\par}\bigskip\bigskip
    {\large\bfseries\mysubtitle{}}\bigskip\par
    \vfill
    %\includegraphics[width=0.5\textwidth]{files/angelfish5.png}
    \vfill
    %{\Large\bfseries\me{}, PhD}\medskip\par
    %{\large\institute{}}%\medskip\par    
    %{\small\today{}}\medskip\par
    %\mybookaddress{} (\href{\mybookgithub}{Github})
\end{titlepage}
%\restorepagecolor
\newpage
\pagestyle{empty}


% \begin{figure}[h]
%     \centering
%     \includegraphics[scale=0.1]{files/IMG_0785.jpeg}
%     \caption{Two of my angels. They are breeding and laying eggs.}\label{fig:angelfish}
%   \end{figure}
  



%   \vfill

% \begin{tcolorbox}[enhanced]
%     \small
%     This \href{\mybookgithub}{work} © 2025 by \me{} is licensed under CC BY-NC-ND 4.0. To view a copy of this license, visit \url{https://creativecommons.org/licenses/by-nc-nd/4.0/}
% \end{tcolorbox}




\renewcommand{\contentsname}{Contents and Summary}
\tableofcontents


\chapter{Preface}




\mainmatter
%\part{Introduction}\label{part:intro}

%\chapter{Should You Apply?}\label{chap:should}

\part{Research}

\part{Funding}

\part{Teaching} 

\part{Advising and Mentoring}

\part{Services}


\part{Administrative}

\chapter{P-Card}

\section{What can you use your P-card for?}

In general, many things including traveling, purchasing supplies, subscriptions, and so on. Here's a list of things that I have used my P-card for:

\begin{itemize}
  \item Registration for conferences, workshops, and so on (\autoref{sec:travel-purchase})
  \item Airfare for traveling to conferences, workshops, and so on
  \item Supplies for my research, such as computer hardware and software, subscriptions (e.g., IEEE, ACM), Cloud services (e.g., AWS, OpenAI), and so on
\end{itemize}


\subsection{Traveling Purchase}\label{sec:travel-purchase}

When you travel for conferences, workshops, and so on, you can use your P-card to pay for some things, but not everything.  You can use your P-card to pay for you or your students
\begin{itemize}
  \item Registration
  \item Airfare
\end{itemize}


But those are the only things you can use your P-card to pay for. Everything including  hotel, Uber, taxi's, food, car rentals, subways, parking, tolls,  you cannot use your P-card. So you (or your students) have to pay for them first, and then get reimbursed. 

You also don't have to use your P-card for registration or airfare and can use your personal credit card instead and get reimbursed later. For example, in some cases when paying for registration in a foreign country, you may not be able to use your P-card because the system does not accept it. In that case, you can use your personal credit card and get reimbursed later. 

It is also important to note that you can only get reimbursed \emph{after} the trip is over. 


\begin{commentbox}[P-Card]
  I try to use my P-card as much as possible to pay for registration and airfare.  So that I don't have to wait for the trip to be over to get reimbursed. This is especially important for students, because they may not have enough money to pay for the trip first and then get reimbursed later. After those purchases, then I don't use my P-card anymore and use my personal credit card instead.
\end{commentbox}

\section{Ordering your P-card}
To order your P-card, you need to fill out a form and submit it to the P-card office. The form is available on the GMU website. You can also ask your department's administrative assistant for help. %TODO: put in website

When it is ready, the P-card office will send the mail card to your department's mailstop. You can then pick it up from there.

\begin{quote}
MS 4A5.  Specifically,  4400 University Dr. MS 4A5\\
Fairfax, VA 22030
\end{quote}


\section{Opening an Amazon Business Account}
If you buy things from Amazon (e.g., supplies, books, and so on), you will want to open an Amazon Business Account. This is different from your personal Amazon account, e.g., they do not charge you tax for purchases. Also you can use your P-card to pay for things through this account.  

To open an Amazon Business Account, email \url{mailto:purch1@gmu.edu} and say that you have a P-Card and want an Amazon business account.



\chapter{Traveling}
\chapter{Training}

\chapter{MFG}


\chapter{Email, Calendar, and VPN}

\section{Email}
GMU Mail uses Outlook 365 and requires you to use OAuth2/Modern Auth to authenticate your GMU credentials.  

\paragraph{Clients} Webmail, Outlook, and Apple Mail are all supported.  The official mails in Android and iOS are also supported. Thunderbird (e.g., on Linux) is also supported but you need to configure it manually as follows (these settings are obtained by logging in to your GMU webmail and going to Settings > Mail > Sync email):

\begin{itemize}
  \item IMAP
  \begin{itemize}
    \item Server: \texttt{outlook.office365.com}
    \item Port: \texttt{993}
    \item Connection security: SSL/TLS
    \item Authentication Method: Leave blank (or OAuth2/Modern Auth if you have it)
    \item Username: Your GMU email address (e.g., tvn@gmu.edu).  Not sure if this is required, but I always use it.
  \end{itemize}
  \item SMTP 
  \begin{itemize}
    \item Server name: \texttt{smtp-mail.outlook.com}
  \item Port: \texttt{587}
  \item Connection Security: STARTTLS
  \item Authentication Method: OAuth2/Modern Auth
  \item Username: Your GMU email address (e.g., tvn@gmu.edu).  Not sure if this is required, but I always use it.
  \end{itemize}
\end{itemize}

\begin{commentbox}
On Linux I can only get Thunderbird to work.  Other email clients do not support OAuth2 well. 
\end{commentbox}

\section{Calendar}
After setting up your email, you can also set up your calendar.  You can use the same clients as above, e.g., Outlook, Apple Calendar, Thunderbird, and so on.  I don't really remember if I had to do anything special to set it up.



\chapter{Personal Website}

\end{document}





 IMAP server
	

outlook.office365.com

IMAP port
	

993

IMAP encryption
	

SSL/TLS

Authentication Method
	

OAuth2/Modern Auth
transparent170x1px.png 	

POP server name
	

outlook.office365.com

POP port
	

995

POP encryption
	

SSL/TLS

Authentication Method
	

OAuth2/Modern Auth

SMTP server name
	

smtp-mail.outlook.com

SMTP port
	

587

SMTP encryption
	

STARTTLS

Authentication Method
	

OAuth2/Modern Auth